\section{Related Work}
\label{sec:related}

\textbf{Parameter-Efficient Fine-Tuning and Model Merging: } Parameter-Efficient Fine-Tuning (PEFT) has emerged as the dominant methodology for adapting large, pre-trained foundation models to specific tasks without incurring the massive computational and storage costs of full-parameter fine-tuning. Prominent PEFT methods include prompt tuning \cite{lester21prompt}, prefix tuning \cite{liliang21prefix}, adapter modules \cite{houlsby19adapter}, and Low-Rank Adaptation (LoRA) \cite{hu21lora}. Among these, LoRA has gained exceptional popularity due to its lack of additional inference latency and high tuning efficiency. When multiple specialized adapters exist, a common practice is \textit{static model merging}, which combines separate model weights into a single set of parameters. Techniques like Task Arithmetic \cite{ilharco22taskarithmetic}, TIES-Merging \cite{yadav23ties}, DARE \cite{yu23dare}, and Model Soups \cite{wortsman22soups} interpolate weight differences to create a unified multi-task model. However, static merging occurs offline and is fundamentally incapable of adapting to non-stationary, real-time input streams at test-time.

\textbf{Dynamic Model Merging and Test-Time Adaptation: } To overcome the limitations of static merging, dynamic model merging frameworks have been developed to blend specialized model parameters on-the-fly based on test-time inputs. Zero-shot centroid alignment (SPS-ZCA) \cite{zhang24spszca} represents a single-pass approach that aligns input features with task centroids in early layers to route activations. Sample-wise Activation Blending (SABLE) \cite{liang23sable} continuously computes ensembling weights at every layer based on the cosine similarity between intermediate activations and task-specific centroids. While SABLE achieves high accuracy, it is stateless and suffers from extreme layer-to-layer ensembling weight jitter because the activations fluctuate rapidly across depth. To suppress this jitter, ChemMerge modeled ensembling updates as first-order non-equilibrium chemical reaction kinetics governed by ordinary differential equations (ODEs). While ChemMerge represents a significant conceptual advance, its reliance on first-order linear concentrations makes it highly sensitive to routing temperatures and reaction constants. In competitive regimes, ChemMerge fails to stabilize, instead amplifying the very jitter it is intended to resolve.

\textbf{Dynamical Systems and Physical Analogies in Deep Learning: } The intersection of deep learning and continuous physical systems has yielded substantial theoretical and practical breakthroughs. The foundational work on Neural Ordinary Differential Equations (Neural ODEs) \cite{chen2018neuralode} established that residual neural networks (ResNets) can be viewed as discrete Euler discretizations of first-order continuous-time dynamical systems \cite{weinan2017proposal, haber2017stable, lu2018beyond, grathwohl2018ffjord}. Other physical formulations, such as Hopfield networks \cite{hopfield1982neural, krotov2016dense} and non-equilibrium thermodynamics \cite{prigogine1978time, nicolis1977self}, have been used to analyze associative memory and feature convergence. Crucially, almost all existing physics-informed deep learning architectures operate as \textit{first-order} gradient-like flows or concentration dynamics. In contrast, GraviMerge introduces a \textit{second-order} physical dynamical system (Newtonian multi-body gravity with velocity and acceleration integration) specifically designed to govern weight-space blending. By introducing physical inertia directly into activation trajectories, GraviMerge represents a fundamentally new category of physics-informed routing that resolves representational volatility where first-order methods fail.
